\chapter{Заключение Тома V: граница архитектуры и программа рождения материи}

Том V начинался с требования найти родительскую архитектуру, которая до
феноменологии фиксирует носитель, пространство состояний, меру, действие и
физическое чтение. Итог не является полной теорией частиц, но он существенно
сильнее каталога гипотез: построен строгий операторно-топологический каркас
материального дефекта и точно доказана граница, на которой этот каркас ещё
не превращается в динамику рождения материи.

Настоящая глава замораживает результат. Она не открывает новый механизм и
не добавляет параметров. Все утверждения разделены на доказанные,
условные и запрещённые к повторному использованию без нового родительского
вывода.

\section{Что доказано}

\subsection{Первичность перехода получила математический язык}

Локальный переход типизирован соответствием Мориты, композиция переходов
--- внутренним тензорным произведением, замкнутый путь --- голономией, а
устойчивое нарушение композиции --- индексом. Связывающая алгебра
$M_{35}(\mathbb C)$ содержит два неравноранговых угла $20$ и $15$ и один
нормированный след. Их суперразмерность равна
\begin{equation}
 \frac{20-15}{35}=\frac17.
 \label{eq:v5-final-superdimension}
\end{equation}

\subsection{Класс пятнадцать построен без ручного проектора ранга пять}

Стабильное тождество
\begin{equation}
 (20-15)\,3=15
\end{equation}
переведено в коэффициентный класс $K_0$ ранга пятнадцать. Расширение
Тёплица и оператор числа на $\ell^2(\mathbb Z)$ дали каноническую
поляризацию, а обменная Real-форма определила ненулевой класс
\begin{equation}
 \kappa_{15}=15\in KO_6(M_{105}(\mathbb C)_{\mathbb R}).
\end{equation}
Явный Real-унитарий имеет ориентированные winding
\begin{equation}
 (+15,-15)
\end{equation}
и граничные индексы противоположных знаков без ложного удвоения до
тридцати.

\subsection{Неплотность замыкания стала операторной теоремой}

В секторе индекса пятнадцать любой фредгольмов представитель удовлетворяет
\begin{equation}
 \dim\ker T+\dim\operatorname{coker}T\geq15.
\end{equation}
Явный оператор Тёплица насыщает эту границу. Поэтому минимальный
нормированный положительный дефект равен
\begin{equation}
 \delta_{\rm cl}=\frac{15}{105}=\frac17.
 \label{eq:v5-final-closure-deficit}
\end{equation}
В полной Real-паре ориентированные индексы сокращаются, но положительная
неплотность сохраняется как $30/210=1/7$.

\subsection{Форма точечного дефекта и его линия определены}

Единственная полностью $SO(3)$-эквивариантная карта
\begin{equation}
 S^2=SO(3)/SO(2)
 \longrightarrow
 \mathbb{RP}^2=SO(3)/O(2)
\end{equation}
равна $\mathbf n\mapsto[\mathbf n]$. Два её глобальных подъёма имеют
степени $+1$ и $-1$ и порождают хопфовы линии $L/L^*$ с
\begin{equation}
 c_1(L)=+1,
 \qquad
 c_1(L^*)=-1.
\end{equation}
После умножения на коэффициентный проектор ранга пятнадцать получаются
классы $+15/-15$. Моритовская градуировка согласует их с $E/E^*$ и
Real-обменом, но не создаёт новый семейный дублет.

\section{Что строго не доказано}

\subsection{Нулевой сектор остаётся допустимым}

Постоянный проектор $P_0$ имеет $dP_0=0$ и нулевой локальный заряд.
Исторический носитель $\mathbb{RP}^3\times S^1$ содержит только торсионный
класс второй когомологии, который ограничивается нулём на локальную сферу;
у позднего кандидата $S^4$ вторая когомология нулевая. Поэтому глобальный
фон не заставляет локальную сферу иметь $c_1=1$.

Минимальное петлевое действие также не выбирает материю глобально:
\begin{equation}
 \mathcal S_{\rm loop}(I)=0,
 \qquad
 \min_{\operatorname{wind}=15}\mathcal S_{\rm loop}=\frac17.
 \label{eq:v5-final-zero-vs-matter-action}
\end{equation}
Второе число является минимумом внутри фиксированного сектора, а не
абсолютным минимумом по всем секторам.

\subsection{Конечная энергия и радиус не выведены}

Проекторная связность допускает четырёхпроизводную кривизну, а
фермионный детерминант может породить квадратичный и четырёхпроизводный
секторы. Но текущий родитель не фиксирует их общий размерный коэффициент,
щель, относительные юкавские амплитуды и схему перенормировки. Поэтому
баланс Деррика остаётся допустимой формой, но не предсказанием радиуса.

\subsection{Полная мера не выбирает одну ориентацию}

Целочисленная WZW-фаза тривиальна после экспоненцирования. Условный
Pfaffian-знак каждой ориентированной ветви нечётен, но полная обменная
Real-пара даёт общий знак $+1$. Ориентация физической Pfaffian line и
правило одной ветви не выведены.

\section{Философская граница результата}

Пусть $M$ обозначает истории с устойчивыми различимыми структурами, а
$O$ --- истории с внутренними наблюдателями. Если наблюдатель является
физической системой с устойчивой памятью, то естественно
\begin{equation}
 O\subset M.
\end{equation}
При любой мере с ненулевой вероятностью $O$ выполнено
\begin{equation}
 \mathbb P(M\mid O)=1.
 \label{eq:v5-final-observer-conditioning}
\end{equation}
Но эта формула не влечёт $\mathbb P(M)=1$.

Следовательно, материя может не быть обязательной во всех математически
допустимых историях, но она обязательна во всякой истории, способной
описывать себя изнутри. Это наблюдательная неизбежность материи, а не
динамическая теорема её рождения.

Такое различение соответствует слабому антропному принципу и литературе
по вероятностям первого лица; см. работы Картера, Хартла--Средницки и
\path{arXiv:gr-qc/0406104}, \path{arXiv:0704.2630},
\path{arXiv:0906.0042}, \path{arXiv:2602.02667}. Оно не разрешает
вводить меру на историях, считать нас типичными или заменять физический
гессиан философской условностью.

\begin{theorem}[Итог Тома V]
Текущая родительская архитектура строго определяет топологию, Real-класс,
ориентационную пару и минимальный положительный дефицит материального
дефекта после фиксации ненулевого сектора. Она не исключает постоянный
нулевой сектор, не делает его динамически неустойчивым и не выводит
конечные массу и радиус. Поэтому Том V завершён как архитектурное и
классификационное замыкание, но не как физическое замыкание рождения
материи.
\end{theorem}

\section{Почему следующий этап должен начинаться с нейтральной пары}

На замкнутом пространстве естественным объектом является не одиночный
нескомпенсированный заряд, а локальная пара противоположных ориентаций.
Уже построенная Real-пара
\begin{equation}
 (V_{15},V_{-15})
\end{equation}
имеет нулевой суммарный ориентированный заряд и ненулевую локальную
положительную неплотность. До физического чтения её следует называть
\emph{дефектом и сопряжённым дефектом}, а не заранее частицей и
античастицей.

Переход
\begin{equation}
 I\longrightarrow(V_{15},V_{-15})
 \label{eq:v5-final-zero-to-real-pair}
\end{equation}
не может целиком лежать в пространстве обратимых унитарных символов:
winding не меняется непрерывно внутри одной компоненты. Поэтому физическое
событие рождения обязано проходить через потерю обратимости, нулевую моду
или rank-defect. Именно такой переход измеряется спектральным потоком.

Литература по спектральному потоку формализует подсчёт пересечений нуля;
см. \path{arXiv:2501.13524} и вещественное обобщение
\path{arXiv:2607.02380}. Литература по тахионной конденсации и
$K$-теории использует аналогичную операторную картину рождения
дефектов из пары с нулевым суммарным зарядом; см.
\path{arXiv:hep-th/0103184}. Эти работы являются методологическими
образцами, а не доказательством проектной динамики.

\section{Предпосылки следующего физического этапа}

Следующий этап допускается только после регистрации общего класса
конфигураций и одного функционала. Порядок проверок фиксируется заранее.

\subsection*{Предпосылка N0: единое пространство конфигураций}

Нужно построить пространство операторов $\mathfrak X$, содержащее
тождественный вакуум, Real-пару и необратимые промежуточные операторы:
\begin{equation}
 I,\ (V_{15},V_{-15})\in\overline{\mathfrak X},
 \qquad
 \mathfrak X_{\rm sing}\ne\varnothing.
\end{equation}
Нельзя сравнивать два сектора, если они определены разными действиями или
разными граничными задачами.

\subsection*{Предпосылка N1: один родительский функционал}

Функционал должен быть конечен на общем классе, иметь одну меру, одно
соглашение о следе и одинаковые контрчлены во всех секторах. Запрещены
секторно-зависимые вычитания и добавление отрицательной энергии только к
дефектной ветви.

\subsection*{Предпосылка N2: физический гессиан нулевого сектора}

После фактора по калибровочным, BV/BRST- и junk-направлениям необходимо
вычислить
\begin{equation}
 \Hess_{I}\mathcal F.
\end{equation}
Рождение материи возможно только если существует физическая отрицательная
мода, ведущая к Real-паре. Если гессиан неотрицателен, текущий родитель
предпочитает пустой вакуум и ветка останавливается до изменения
архитектуры.

\subsection*{Предпосылка N3: нелинейное насыщение}

Отрицательная мода недостаточна. Должна существовать конечная стационарная
конфигурация с локальными классами $+15/-15$, положительным физическим
гессианом после удаления симметрий и без ручной фиксации топологической
границы.

\subsection*{Предпосылка N4: пространственная локализация}

Один функционал должен дать конечные энергию, радиус и расстояние между
сопряжёнными дефектами. Баланс типа
\begin{equation}
 E(R)=aR+\frac bR
\end{equation}
принимается только если $a$ и $b$ следуют из одного оператора и одной
нормировки.

\subsection*{Предпосылка N5: квантовая мера}

Требуется построить фермионный оператор, determinant/Pfaffian line,
проверить аномалии, статистику и устойчивость результата после
квантовых поправок. Нельзя использовать модуль полной Real-пары в энергии
и одну её половину только в фазе.

\subsection*{Предпосылка N6: космологическое рождение и плотность}

Только после вывода фазового перехода допустимо применять механизм
Киббла--Зурека к плотности рождающихся дефектов. Этот механизм связывает
скорость прохождения критической точки с масштабом корреляции и плотностью
дефектов; см. \path{arXiv:2402.18533} и
\path{arXiv:2602.08759}. Он не создаёт отсутствующий фазовый переход и не
фиксирует его параметры.

\subsection*{Предпосылка N7: физическое чтение}

Названия ``частица'', ``античастица'', ``масса'', ``заряд'' и
``поколение'' разрешаются только после построения локальных наблюдаемых и
их преобразований относительно группы Стандартной модели. Затем
регистрируются безразмерные слепые проверки; массовые данные не могут
выбирать действие задним числом.

\section{Критерии немедленной остановки}

Следующий этап закрывается без феноменологии, если выполняется хотя бы один
пункт:
\begin{enumerate}
 \item общий путь от $I$ к Real-паре требует ручной смены пространства
 состояний;
 \item физический гессиан нуля не имеет отрицательной моды;
 \item стабилизация требует независимых коэффициентов $E_2$ и $E_4$;
 \item конечный радиус существует только после выбора внешнего масштаба,
 не являющегося разрешённым обучающим входом;
 \item квантовая мера обнуляет или дестабилизирует локальный класс;
 \item физическое чтение требует новых состояний вне доказанного бюджета.
\end{enumerate}

\section{Итоговая программа}

Следующий этап не должен снова спрашивать, какая линия или топология
описывает уже существующую частицу. Его первый вопрос единственен:
\begin{equation}
 \boxed{
 \begin{gathered}
 \text{может ли нулевой вакуум одного родительского действия}\\
 \text{сам породить устойчивую нейтральную Real-пару дефектов?}
 \end{gathered}}
 \label{eq:v5-final-next-question}
\end{equation}

Положительный ответ начинает физику рождения материи. Отрицательный ответ
оставляет строгий дефектный каркас как классификацию возможной
наблюдаемой ветви и фиксирует наблюдательную, но не динамическую
неизбежность материи.

\begin{protocol}[Финальная заморозка Тома V]
\begin{enumerate}
 \item родительский моритовский каркас: \textbf{построен};
 \item обменный Real-класс: \textbf{$15$};
 \item минимальный положительный дефицит: \textbf{$1/7$};
 \item эквивариантный хопфов подъём: \textbf{$c_1=\pm1$};
 \item постоянный нулевой сектор: \textbf{допустим};
 \item глобальный топологический выбор материи: \textbf{не получен};
 \item наблюдательная условность $\mathbb P(M\mid O)=1$:
 \textbf{сформулирована};
 \item конечные масса и радиус: \textbf{не выведены};
 \item физическая мера: \textbf{не замкнута};
 \item Том V как архитектурный аудит: \textbf{завершён};
 \item первый вопрос следующего этапа:
 \textbf{неустойчивость нуля к нейтральной Real-паре}.
\end{enumerate}
\end{protocol}